% from ega4-21b-fr.tex
\oldpage[IV-4]{294}
qui est donc un cycle $1$-codimensionnel sur $X'$ à coefficients \emph{rationnels}.
On a encore défini ainsi un homomorphisme $f^*$ de groupes ordonnés, vérifiant \sref{IV.21.10.3.1}, et par suite un homomorphisme de faisceaux de groupes commutatifs
\[
  \psi^*(\Gamma_T(\mathscr{Z}^1_X))\to\mathscr{Z}^1_{X'}\otimes_{\bb{Z}}\bb{Q}.
\]

Lorsque $f$ vérifie les conditions précédentes pour \emph{toute partie fermée $T$ rare} dans $X$, c'est-à-dire que pour tout $x'\in X'^{(1)}$, ou bien $x=f(x')$ est point maximal de $X$, ou bien $x'$ vérifie l'une des conditions (ii), (iii) ou (iv), $f^*(Z)$ est alors défini pour tout cycle $1$-codimensionnel $Z$ sur $X$, et on a défini un homomorphisme de faisceaux de groupes commutatifs ordonnés
\[
  \psi^*(\mathscr{Z}^1_X)\to\mathscr{Z}^1_{X'}\otimes_{\bb{Z}}\bb{Q},
\]
d'où, par tensorisation, un homomorphisme de faisceaux de $\bb{Q}$-espaces vectoriels ordonnés
\[
\label{IV.21.10.11.1}
  \psi^*(\mathscr{Z}^1_X\otimes_{\bb{Z}}\bb{Q})\to\mathscr{Z}^1_{X'}\otimes_{\bb{Z}}\bb{Q}.
  \tag{21.10.11.1}
\]
\end{env}

\begin{remark}[21.10.12]
\label{IV.21.10.12}
Lorsqu'on est dans la situation de \sref{IV.21.10.10}, on peut effectivement, pour des cycles $1$-codimensionnels $Z$ sur $X$, avoir pour $f^*(Z)$ des cycles $1$-codimensionnels à coefficients \emph{non entiers}; autrement dit les nombres $\mu_{x'}$ peuvent être non entiers.
On en a un exemple en prenant l'anneau intègre complet $A$ de $(\mathbf{0}_{\mathrm{IV}},6.15.11)$, (ii) et sa clôture intégrale $A'$: le point fermé $x'$ de $\Spec(A')$ vérifie la condition (iv) de \sref{IV.21.10.10} et l'on a $\mu_{x'}=\frac12$.
\end{remark}

\begin{lemma}[21.10.13]
\label{IV.21.10.13}
Soient $A$ un anneau local noethérien de dimension $1$, $t$ un élément régulier de $A$ appartenant à l'idéal maximal $\mathfrak{m}$ (ce qui implique que $\mathfrak{m}\notin\operatorname{Ass}(A)$).
\begin{enumerate}[label=(\roman*)]
  \item Pour tout $A$-module $M$ de type fini, le noyau $N_t(M)$ et le conoyau $P_t(M)$ de l'homothétie $t_M:M\to M$ de rapport $t$ sont de longueurs finies. On pose $d_t(M)=\operatorname{long}P_t(M)-\operatorname{long}N_t(M)$.
  \item Si $0\to M'\to M\to M''\to0$ est une suite exacte de $A$-modules de type fini, on a $d_t(M)=d_t(M')+d_t(M'')$.
  \item On a $d_t(M)\geq0$ pour tout $A$-module $M$ de type fini; pour que $d_t(M)=0$, il faut et il suffit que $M$ soit de longueur finie.
  \item Si $K$ est l'anneau total des fractions de $A$ et si $M\otimes_AK$ est un $K$-module libre de rang $n$, on a $d_t(M)=n\cdot d_t(A)=n\cdot\operatorname{long}(A/tA)$.
  \item Si $M$ vérifie les hypothèses de (iv) et si de plus $t$ est $M$-régulier, on a $\operatorname{long}(M/tM)=n\cdot\operatorname{long}(A/tA)$.
\end{enumerate}
\end{lemma}

\begin{proof}
(i) $\Spec(A)$ est formé du point $\mathfrak{m}$ et des idéaux premiers minimaux $\mathfrak{p}_i$;
comme par hypothèse $t\notin\mathfrak{p}_i$ pour tout $i$ (Bourbaki, \emph{Alg. comm.}, chap.~IV, \textsection~1, n\textsuperscript{o}~1, cor.~3 de la prop.~2), l'image de $t$ dans chacun des $A_{\mathfrak{p}_i}$ est inversible, et les supports des $A$-modules de type fini $N_t(M)$ et $P_t(M)$ sont donc vides ou réduits à $\mathfrak{m}$;
on en conclut $(\mathbf{0},16.1.10)$ que ces modules sont de longueur finie.

(ii) Puisque $t$ est régulier, on a une suite exacte
\[
  0\to A\xrightarrow{t_A}A\to A/tA\to0
\]
